\documentclass[11pt]{article}
\usepackage[a4paper,margin=22mm]{geometry}
\usepackage{fontspec}
\defaultfontfeatures{Ligatures=TeX}
\setmainfont{Times New Roman}
\setsansfont{Times New Roman}
\setmonofont{Courier New}
\usepackage{amsmath,mathtools,unicode-math}
\setmathfont{STIX Two Math}
\usepackage{graphicx,booktabs,tabularx,xltabular,array,ragged2e}
\usepackage{caption,float,placeins,needspace,setspace,microtype}
\usepackage{xcolor}
\usepackage[hidelinks]{hyperref}
\hypersetup{pdftitle={Supplementary Material: Scale-Explicit and Uncertainty-Qualified Assessment of Tangent-Only Matrix Perturbations in Coupled Finite-Element Linearizations}}
\renewcommand{\tabularxcolumn}[1]{>{\RaggedRight\arraybackslash}p{#1}}
\renewcommand{\arraystretch}{1.14}
\captionsetup{font=small,labelfont=bf,skip=5pt}
\setlength{\parindent}{1.5em}
\setlength{\parskip}{0.08em}
\setlength{\emergencystretch}{2em}
\setstretch{1.04}
\title{\bfseries Supplementary Material\\[0.4em]\large Scale-Explicit and Uncertainty-Qualified Assessment of Tangent-Only Matrix Perturbations in Coupled Finite-Element Linearizations}
\author{Bing Liang \and Yangqi Ma \and Weiji Sun\thanks{Corresponding author: sunweiji-1231@163.com} \and Shi He}
\date{}

\begin{document}
\maketitle
\renewcommand{\thetable}{S\arabic{table}}
\renewcommand{\thefigure}{S\arabic{figure}}

\section*{Purpose and scope}
This supplement preserves exact numerical values, uncertainty components, precision and identity checks, tooling provenance, and the complete claim--evidence matrix supporting the main manuscript. M8 and M16 are independent frozen finite-element linearizations rather than a mesh or refinement sequence. The material below documents the existing calculations; no new scientific case was executed for the public-release v0.6 revision.

\section{Synthetic analytical verification}
\begin{table}[htbp]\centering\small
\caption{Pre-specified reproduction of the P2 directional construction.}\label{tab:S1}
\begin{tabularx}{0.98\textwidth}{@{}>{\raggedright\arraybackslash}p{0.07\textwidth}XXXXX@{}}
\toprule
$\varepsilon$ & $\mu_J$ (matched) & $\rho_1$ (matched) & $\rho_g$ (matched) & aligned $\eta_z$ & orthogonal $\eta_z$ \\
\midrule
1e-12 & 1.0000e-12 & 1.0000e-12 & 1.0000e-12 & 1.0000e-12 & 0 \\
1e-10 & 1.0000e-10 & 1.0000e-10 & 1.0000e-10 & 1.0000e-10 & 0 \\
1e-8 & 1.0000e-08 & 1.0000e-08 & 1.0000e-08 & 1.0000e-08 & 0 \\
1e-6 & 1.0000e-06 & 1.0000e-06 & 1.0000e-06 & 1.0000e-06 & 0 \\
1e-4 & 1.0000e-04 & 1.0000e-04 & 9.9990e-05 & 9.9990e-05 & 0 \\
1e-2 & 1.0000e-02 & 1.0000e-02 & 9.9010e-03 & 9.9010e-03 & 0 \\
\bottomrule
\end{tabularx}

\end{table}

\begin{table}[htbp]\centering\small
\caption{Pre-specified P3 coordinate-transformation and redeclaration calculation.}\label{tab:S2}
\begin{tabularx}{0.98\textwidth}{@{}>{\raggedright\arraybackslash}p{0.07\textwidth}XXXXX@{}}
\toprule
$s$ & transformed $\widehat{\Delta J}_{22}$ & transformed physical $(\Delta J)_{22}$ & transformed $\eta_z$ & redeclared physical $(\Delta J)_{22}$ & redeclared $\eta_z$ \\
\midrule
1 & 1.0000e-06 & 1.0000e-06 & 1.0000e-06 & 1.0000e-06 & 1.0000e-06 \\
10 & 1.0000e-05 & 1.0000e-06 & 1.0000e-06 & 1.0000e-07 & 1.0000e-07 \\
1e3 & 1.0000e-03 & 1.0000e-06 & 1.0000e-06 & 1.0000e-09 & 1.0000e-09 \\
1e6 & 1.0000e+00 & 1.0000e-06 & 1.0000e-06 & 1.0000e-12 & 1.0000e-12 \\
\bottomrule
\end{tabularx}

\end{table}

\section{Frozen operand provenance}

\begin{tabularx}{\linewidth}{@{}p{0.12\linewidth}p{0.12\linewidth}p{0.24\linewidth}X@{}}
\toprule
\textbf{Case} & \textbf{Dimension} & \textbf{Snapshot SHA-256} & \textbf{Declared operand contract} \\
\midrule
M8 & 196 & \texttt{ACF00D823AAE\allowbreak{}8BC073F14ECF\allowbreak{}A322AED1B004\allowbreak{}CBBA85D5FA83\allowbreak{}416FEC4D24D8\allowbreak{}7482} & Frozen IEEE binary64 $J$, $J_g$, $b$, free map, and scales; transformed free coordinates; Euclidean norm. \\
M16 & 900 & \texttt{22A51CA5BB44\allowbreak{}03D1EB9499F4\allowbreak{}588D6EF19CE7\allowbreak{}D6C3636429D5\allowbreak{}EE2FCA6AE3AB\allowbreak{}207E} & Independently frozen IEEE binary64 $J$, $J_g$, $b$, free map, and scales; transformed free coordinates; Euclidean norm. \\
\bottomrule
\end{tabularx}


The M8 and M16 matrix snapshots were each reproduced bitwise by an independent capture witness. Each offline application reconstructed every declared binary64 operand exactly before high-precision comparison. The perturbation identity $J_g=J+\Delta J$ was checked elementwise before either response route was classified.

\section{Exact route-specific results}
\footnotesize
\setlength{\tabcolsep}{3.5pt}

\begin{xltabular}{0.99\textwidth}{@{}p{0.08\textwidth}p{0.18\textwidth}p{0.28\textwidth}X@{}}
\caption{Exact M8 and M16 route-specific values copied from the registered decision files.}\label{tab:S4}\\
\toprule
\textbf{Case} & \textbf{Route} & \textbf{Quantity} & \textbf{Exact registered value} \\
\midrule\endfirsthead
\toprule
\textbf{Case} & \textbf{Route} & \textbf{Quantity} & \textbf{Exact registered value} \\
\midrule\endhead
M8 & Direct difference & Observed response & \texttt{0.0000000000\allowbreak{}029185014536\allowbreak{}758156190032\allowbreak{}956007920590\allowbreak{}207248424420\allowbreak{}081982240251\allowbreak{}520232804179\allowbreak{}861608795259\allowbreak{}826364819968\allowbreak{}05201} \\
M8 & Direct difference & Uncertainty bound & \texttt{0.0000000000\allowbreak{}005880137782\allowbreak{}404462163207\allowbreak{}715781856199\allowbreak{}927166001714\allowbreak{}634199187384\allowbreak{}669088463464\allowbreak{}446774100523\allowbreak{}218868676802\allowbreak{}748951} \\
M8 & Direct difference & Lower interval bound & \texttt{0.0000000000\allowbreak{}023304876754\allowbreak{}353694026825\allowbreak{}240226064390\allowbreak{}280082422705\allowbreak{}447783052866\allowbreak{}851144340715\allowbreak{}414834694736\allowbreak{}607496143165\allowbreak{}30306} \\
M8 & Direct difference & Upper interval bound & \texttt{0.0000000000\allowbreak{}035065152319\allowbreak{}162618353240\allowbreak{}671789776790\allowbreak{}134414426134\allowbreak{}716181427636\allowbreak{}189321267644\allowbreak{}308382895783\allowbreak{}045233496770\allowbreak{}80096} \\
M8 & Direct difference & Guaranteed relative-error upper bound & \texttt{0.2523136184\allowbreak{}921452489334\allowbreak{}362403820407\allowbreak{}858086971785\allowbreak{}578796069759\allowbreak{}756449422942\allowbreak{}232194245526\allowbreak{}536411060357\allowbreak{}365522} \\
M8 & Direct difference & Decision & \texttt{DISTINGUISHABLE\_FROM\_ZERO} \\
M8 & Response equation & Observed response & \texttt{0.0000000000\allowbreak{}029200142254\allowbreak{}489211942473\allowbreak{}568021747142\allowbreak{}861853194538\allowbreak{}287404086512\allowbreak{}718756577010\allowbreak{}642403187437\allowbreak{}023694400712\allowbreak{}51828} \\
M8 & Response equation & Uncertainty bound & \texttt{0.0000000000\allowbreak{}000000000002\allowbreak{}588299567178\allowbreak{}916551202655\allowbreak{}296880895539\allowbreak{}935734612342\allowbreak{}619760039891\allowbreak{}099523264193\allowbreak{}484210461160\allowbreak{}608761377537\allowbreak{}887} \\
M8 & Response equation & Lower interval bound & \texttt{0.0000000000\allowbreak{}029200142251\allowbreak{}900912375294\allowbreak{}651470544487\allowbreak{}564972298998\allowbreak{}351669474170\allowbreak{}098996537119\allowbreak{}542879923243\allowbreak{}539483939551\allowbreak{}90952} \\
M8 & Response equation & Upper interval bound & \texttt{0.0000000000\allowbreak{}029200142257\allowbreak{}077511509652\allowbreak{}484572949798\allowbreak{}158734090078\allowbreak{}223138698855\allowbreak{}338516616901\allowbreak{}741926451630\allowbreak{}507904861873\allowbreak{}12704} \\
M8 & Response equation & Guaranteed relative-error upper bound & \texttt{0.0000000000\allowbreak{}886399643142\allowbreak{}293166771151\allowbreak{}775600468128\allowbreak{}269760609861\allowbreak{}951448626546\allowbreak{}399600511363\allowbreak{}704411828880\allowbreak{}102699434908\allowbreak{}2996} \\
M8 & Response equation & Decision & \texttt{DISTINGUISHABLE\_FROM\_ZERO} \\
M16 & Direct difference & Observed response & \texttt{0.0000000000\allowbreak{}246337644077\allowbreak{}072171041837\allowbreak{}867972547031\allowbreak{}193063729610\allowbreak{}422125441092\allowbreak{}925581517718\allowbreak{}827382698855\allowbreak{}426123142196\allowbreak{}727} \\
M16 & Direct difference & Uncertainty bound & \texttt{0.0000000000\allowbreak{}039707999109\allowbreak{}429423245734\allowbreak{}513521266219\allowbreak{}270598709891\allowbreak{}319683247093\allowbreak{}224268735595\allowbreak{}418035162208\allowbreak{}005533422405\allowbreak{}12158} \\
M16 & Direct difference & Lower interval bound & \texttt{0.0000000000\allowbreak{}206629644967\allowbreak{}642747796103\allowbreak{}354451280811\allowbreak{}922465019719\allowbreak{}102442193999\allowbreak{}701312782123\allowbreak{}409347536647\allowbreak{}420589719791\allowbreak{}6054} \\
M16 & Direct difference & Upper interval bound & \texttt{0.0000000000\allowbreak{}286045643186\allowbreak{}501594287572\allowbreak{}381493813250\allowbreak{}463662439501\allowbreak{}741808688186\allowbreak{}149850253314\allowbreak{}245417861063\allowbreak{}431656564601\allowbreak{}8485} \\
M16 & Direct difference & Guaranteed relative-error upper bound & \texttt{0.1921699043\allowbreak{}505955464754\allowbreak{}606186404427\allowbreak{}103889558168\allowbreak{}489583509620\allowbreak{}597278991055\allowbreak{}977662511355\allowbreak{}172628645464\allowbreak{}869451} \\
M16 & Direct difference & Decision & \texttt{DISTINGUISHABLE\_FROM\_ZERO} \\
M16 & Response equation & Observed response & \texttt{0.0000000000\allowbreak{}246367184609\allowbreak{}706229802356\allowbreak{}176545021502\allowbreak{}495517486557\allowbreak{}655708640777\allowbreak{}051756475467\allowbreak{}152072508356\allowbreak{}411351653360\allowbreak{}519} \\
M16 & Response equation & Uncertainty bound & \texttt{0.0000000000\allowbreak{}000000000069\allowbreak{}304826171673\allowbreak{}561801211534\allowbreak{}161663831658\allowbreak{}382676666413\allowbreak{}069558694185\allowbreak{}951345727480\allowbreak{}057886539594\allowbreak{}297019202223} \\
M16 & Response equation & Lower interval bound & \texttt{0.0000000000\allowbreak{}246367184540\allowbreak{}401403630682\allowbreak{}614743809968\allowbreak{}333853654899\allowbreak{}273031974363\allowbreak{}982197781281\allowbreak{}200726780876\allowbreak{}353465113766\allowbreak{}222} \\
M16 & Response equation & Upper interval bound & \texttt{0.0000000000\allowbreak{}246367184679\allowbreak{}011055974029\allowbreak{}738346233036\allowbreak{}657181318216\allowbreak{}038385307190\allowbreak{}121315169653\allowbreak{}103418235836\allowbreak{}469238192954\allowbreak{}816} \\
M16 & Response equation & Guaranteed relative-error upper bound & \texttt{0.0000000002\allowbreak{}813070511032\allowbreak{}623424952586\allowbreak{}212313222893\allowbreak{}810817729228\allowbreak{}669242584304\allowbreak{}341140586873\allowbreak{}408936987001\allowbreak{}183002654740\allowbreak{}493} \\
M16 & Response equation & Decision & \texttt{DISTINGUISHABLE\_FROM\_ZERO} \\
\bottomrule
\end{xltabular}

\normalsize

\section{Precision, residual, identity, and repeatability evidence}
\footnotesize
\setlength{\tabcolsep}{3.2pt}

\begin{xltabular}{0.99\textwidth}{@{}p{0.08\textwidth}p{0.18\textwidth}p{0.22\textwidth}Xp{0.10\textwidth}p{0.08\textwidth}@{}}
\caption{Independent-reference, identity, and repeatability gates.}\label{tab:S5}\\
\toprule
\textbf{Case} & \textbf{Audit} & \textbf{Object} & \textbf{Maximum registered difference} & \textbf{Gate} & \textbf{Status} \\
\midrule\endfirsthead
\toprule
\textbf{Case} & \textbf{Audit} & \textbf{Object} & \textbf{Maximum registered difference} & \textbf{Gate} & \textbf{Status} \\
\midrule\endhead
M8 & Precision ladder & 120/180 dps & \texttt{8.8057092380\allowbreak{}884434187176\allowbreak{}962461947943\allowbreak{}092333470546\allowbreak{}718574519151\allowbreak{}728781821208\allowbreak{}054678586278\allowbreak{}823203520491\allowbreak{}55351e-121} & $\leq 10^{-80}$ & PASS \\ 
M8 & Reference residual & Parent and response systems & \texttt{2.5509756331\allowbreak{}440603586511\allowbreak{}214959945542\allowbreak{}624542125528\allowbreak{}335796688014\allowbreak{}351623728996\allowbreak{}256868809424\allowbreak{}330779651386\allowbreak{}98482e-124} & $\leq 10^{-90}$ & PASS \\ 
M8 & Response identity & Parent difference versus response solve & \texttt{9.1849085476\allowbreak{}145843658430\allowbreak{}984457176152\allowbreak{}739279490697\allowbreak{}844155754981\allowbreak{}011387453725\allowbreak{}230772676985\allowbreak{}998487572365\allowbreak{}00405e-123} & $\leq 10^{-80}$ & PASS \\ 
M8 & Repeatability & 120/180 dps independent process & \texttt{0; 0} & $\leq 10^{-80}$ & PASS \\ 
M16 & Precision ladder & 120/180 dps & \texttt{3.9700701327\allowbreak{}504386317003\allowbreak{}831479435048\allowbreak{}947013122943\allowbreak{}519143716924\allowbreak{}264139187506\allowbreak{}177819881847\allowbreak{}658177636420\allowbreak{}02851e-121} & $\leq 10^{-80}$ & PASS \\ 
M16 & Reference residual & Parent and response systems & \texttt{4.1440816367\allowbreak{}892560429801\allowbreak{}527838154080\allowbreak{}514775758958\allowbreak{}038539094508\allowbreak{}388888254696\allowbreak{}147598977293\allowbreak{}450991336676\allowbreak{}71501e-124} & $\leq 10^{-90}$ & PASS \\ 
M16 & Response identity & Parent difference versus response solve & \texttt{3.0422633711\allowbreak{}275487636633\allowbreak{}980676903715\allowbreak{}062449817277\allowbreak{}763274882660\allowbreak{}901630223123\allowbreak{}725003858668\allowbreak{}546771128132\allowbreak{}64999e-122} & $\leq 10^{-80}$ & PASS \\ 
M16 & Repeatability & 120/180 dps independent process & \texttt{0; 0} & $\leq 10^{-80}$ & PASS \\ 
\bottomrule
\end{xltabular}

\normalsize

The reference precision ladder used 120 and 180 decimal digits for each application. M8 consumed exactly 467.59207850007806~s with recorded peak memory of 115154944 bytes. M16 consumed exactly 37512.160941799986~s with recorded peak memory of 1452195840 bytes. These values document two offline applications; they are not a complexity fit or an online-estimator benchmark.

\section{Route-specific uncertainty components}
\footnotesize
\setlength{\tabcolsep}{3.2pt}

\begin{xltabular}{0.99\textwidth}{@{}p{0.08\textwidth}p{0.19\textwidth}p{0.20\textwidth}Xp{0.10\textwidth}@{}}
\caption{Route-specific uncertainty-component coverage.}\label{tab:S6}\\
\toprule
\textbf{Case} & \textbf{Route} & \textbf{Component} & \textbf{Registered component bound} & \textbf{Status} \\
\midrule\endfirsthead
\toprule
\textbf{Case} & \textbf{Route} & \textbf{Component} & \textbf{Registered component bound} & \textbf{Status} \\
\midrule\endhead
M8 & Direct difference & Parent-$z$ & \texttt{0.0000000000\allowbreak{}003226016562\allowbreak{}676377291832\allowbreak{}930729419859\allowbreak{}852938628308\allowbreak{}420601477443\allowbreak{}639054955629\allowbreak{}477059647754\allowbreak{}299784885800\allowbreak{}860621} & COVERED \\ 
M8 & Direct difference & Parent-$z_g$ & \texttt{0.0000000000\allowbreak{}002654033664\allowbreak{}645770837841\allowbreak{}684095747019\allowbreak{}844356589723\allowbreak{}578280133643\allowbreak{}990225258788\allowbreak{}469947115176\allowbreak{}298525210742\allowbreak{}069229} & COVERED \\ 
M8 & Direct difference & Final subtraction & \texttt{0.0000000000\allowbreak{}000000087555\allowbreak{}082314033533\allowbreak{}100956689320\allowbreak{}229870783682\allowbreak{}635317576297\allowbreak{}039808249046\allowbreak{}499767337592\allowbreak{}620558580259\allowbreak{}81910067742} & COVERED \\ 
M8 & Response equation & Forcing formation & \texttt{0.0} & COVERED \\ 
M8 & Response equation & Inherited parent & \texttt{0.0000000000\allowbreak{}000000000000\allowbreak{}000011461113\allowbreak{}170257339148\allowbreak{}722517767779\allowbreak{}506858043093\allowbreak{}261374428512\allowbreak{}921636009578\allowbreak{}312627822241\allowbreak{}702445839755\allowbreak{}60054732} & COVERED \\ 
M8 & Response equation & Response solve & \texttt{0.0000000000\allowbreak{}000000000000\allowbreak{}000180735804\allowbreak{}346278039810\allowbreak{}310987544425\allowbreak{}429697255409\allowbreak{}874964364130\allowbreak{}030649103282\allowbreak{}552895454713\allowbreak{}816921970032\allowbreak{}0866423} & COVERED \\ 
M16 & Direct difference & Parent-$z$ & \texttt{0.0000000000\allowbreak{}019983166032\allowbreak{}948826973316\allowbreak{}521370735139\allowbreak{}410993713261\allowbreak{}884266407930\allowbreak{}866507754795\allowbreak{}342051914738\allowbreak{}870258383077\allowbreak{}58855} & COVERED \\ 
M16 & Direct difference & Parent-$z_g$ & \texttt{0.0000000000\allowbreak{}019724636210\allowbreak{}036355150009\allowbreak{}999633696356\allowbreak{}207605901402\allowbreak{}088592157309\allowbreak{}632550972744\allowbreak{}979627856778\allowbreak{}648345399536\allowbreak{}72571} & COVERED \\ 
M16 & Direct difference & Final subtraction & \texttt{0.0000000000\allowbreak{}000000196866\allowbreak{}444241122407\allowbreak{}992516834723\allowbreak{}651999095227\allowbreak{}346824681852\allowbreak{}725210008055\allowbreak{}096355390690\allowbreak{}486929639790\allowbreak{}8073150567} & COVERED \\ 
M16 & Response equation & Forcing formation & \texttt{0.0} & COVERED \\ 
M16 & Response equation & Inherited parent & \texttt{0.0000000000\allowbreak{}000000000000\allowbreak{}000070994467\allowbreak{}310998839243\allowbreak{}000167338717\allowbreak{}927417923594\allowbreak{}866111139395\allowbreak{}779853017226\allowbreak{}895922149462\allowbreak{}375011304910\allowbreak{}81933567} & COVERED \\ 
M16 & Response equation & Response solve & \texttt{0.0000000000\allowbreak{}000000000000\allowbreak{}004963781817\allowbreak{}969057629693\allowbreak{}103945555487\allowbreak{}579252590699\allowbreak{}178538847648\allowbreak{}306411221536\allowbreak{}674972641507\allowbreak{}398272667592\allowbreak{}811359} & COVERED \\ 
\bottomrule
\end{xltabular}

\normalsize

All declared uncertainty components covered their independently evaluated reference errors in the two frozen applications. Coverage is route- and case-specific and does not certify either arithmetic route for other operators.

\section{Tooling-attempt provenance}
An initial G2-11 attempt terminated before scientific evaluation because a NumPy Boolean value could not be serialized to JSON. That failed attempt was retained. The only correction converted one NumPy Boolean result to the built-in Python Boolean type. The complete execution then restarted from its first row rather than continuing from the stopping point. Neither attempt invoked Abaqus, a user element, finite-element assembly, a Newton replay, or a new mesh.

\section{Complete claim--evidence matrix}
The CSV file \texttt{Paper2\_claim\_evidence\_matrix\_v0.4.csv} is the authoritative record. Table~\ref{tab:S7} is a typeset presentation view retained in portrait orientation.

\scriptsize
\setlength{\tabcolsep}{2.1pt}
\begin{xltabular}{0.99\linewidth}{@{}p{0.07\linewidth}p{0.12\linewidth}p{0.27\linewidth}p{0.22\linewidth}X@{}}
\caption{Expanded v0.4 claim--evidence matrix. The CSV is authoritative.}\label{tab:S7}\\
\toprule
\textbf{ID} & \textbf{Tier} & \textbf{Claim} & \textbf{Allowed strength} & \textbf{Prohibited extrapolation} \\
\midrule\endfirsthead
\toprule
\textbf{ID} & \textbf{Tier} & \textbf{Claim} & \textbf{Allowed strength} & \textbf{Prohibited extrapolation} \\
\midrule\endhead
C01 & T1\_\allowbreak{}THEOREM & Matrix-space smallness alone is insufficient to determine a realized directional solution response. & State as an exact insufficiency result under declared norms and probes. & A small guard is generally unsafe or harmless. \\
C02 & T1\_\allowbreak{}THEOREM & For common coordinates and right-hand side, the finite response satisfies the exact guarded-inverse identity. & Use proves or satisfies under the assumptions. & Changed residual roots or nonlinear trajectories follow the same identity. \\
C03 & T1\_\allowbreak{}THEOREM & The realized response is bounded by the exact guarded-inverse amplification in a compatible subordinate norm. & Use bounds under the declared norm. & rho\_\allowbreak{}g is the observed response. \\
C04 & T1\_\allowbreak{}THEOREM & When rho\_\allowbreak{}1 is below one, the finite response admits the Neumann-series bound with its denominator retained. & Use finite bound with assumptions. & Drop the denominator outside the asymptotic regime. \\
C05 & T1\_\allowbreak{}THEOREM & kappa(J)mu\_\allowbreak{}J is a worst-case normwise comparison and not a directional mechanism classifier. & Use as a loose worst-case comparison. & Condition number alone identifies the observed mechanism. \\
C06 & T1\_\allowbreak{}THEOREM & Matched matrix amplitude and worst-direction metrics can yield different realized responses for one fixed probe. & Use proves non-determination. & Typicality, frequency, mesh law, safety, or harm. \\
C07 & T2\_\allowbreak{}VALIDATED\_\allowbreak{}REGISTERED & The registered P2 calculation reproduces the exact metric match and response distinction for six frozen perturbation magnitudes. & Use validates or reproduces. & General operator-family frequency. \\
C08 & T1\_\allowbreak{}THEOREM & A physical-coordinate identity guard generally transforms to epsilon D\_\allowbreak{}r\textasciicircum{}(-1)D\_\allowbreak{}x rather than the same scalar identity. & Use proves coordinate dependence of the representation. & Every physical identity guard is harmful. \\
C09 & T1\_\allowbreak{}THEOREM & A dimensionless identity guard maps to a generally field-dependent physical perturbation. & Use proves the reverse mapping. & The selected scaling is optimal or uniquely physical. \\
C10 & T2\_\allowbreak{}VALIDATED\_\allowbreak{}REGISTERED & Consistent transformation preserves the registered physical operator and response across the tested coordinate descriptions. & Use validates equivalent transformation in the tested audit. & Universal invariance of independently redeclared guards. \\
C11 & T2\_\allowbreak{}VALIDATED\_\allowbreak{}REGISTERED & Redeclaring the same numeric identity after scaling defines a different physical perturbation for every registered scale factor above one. & Use demonstrates within the frozen audit. & One coordinate system is safer or optimal. \\
C12 & T1\_\allowbreak{}THEOREM & Direct subtraction supports an observability statement only when parent-solution and subtraction-error bounds cover the declared truth lane. & Use conditional on the registered error model. & Direct subtraction is universally unreliable or always reliable. \\
C13 & T2\_\allowbreak{}VALIDATED\_\allowbreak{}REGISTERED & Source-truth and binary-operand truth require separate, non-mixed uncertainty lanes in the registered P4-r1 audit. & Use defines the validated two-truth procedure. & One truth lane can be substituted for the other. \\
C14 & T2\_\allowbreak{}VALIDATED\_\allowbreak{}REGISTERED & The historical direct uncertainty model is reproduced exactly at 287/300 coverage with three severe undercoverage cases. & Use reproduces and discloses. & The guard caused the undercoverage. \\
C15 & T2\_\allowbreak{}VALIDATED\_\allowbreak{}REGISTERED & The revised source and binary envelopes cover all 300 registered development cases with no severe, component, subtraction, or nonfinite failures. & Use validates coverage on the registered development set. & Universal error bound or classification sharpness. \\
C16 & T2\_\allowbreak{}VALIDATED\_\allowbreak{}REGISTERED & The revised source and binary envelopes cover all 300 disjoint registered holdout cases with no severe, component, subtraction, or nonfinite failures. & Use validates coverage on the registered holdout. & Universal coverage or WP27 class validation. \\
C17 & T4\_\allowbreak{}LIMITATION & Coverage validation does not qualify bound tightness or observability-class sharpness. & Use explicitly as a limitation. & Tight, efficient, sharp, or calibrated class thresholds. \\
C18 & T3\_\allowbreak{}BOUNDED\_\allowbreak{}DIAGNOSTIC & A response solve can reduce subtraction error in selected small-guard cases, but it is not an authoritative truth route without its own validated uncertainty bound. & Use motivates an independent response lane. & Stable lane is exact truth or always better. \\
C19 & T3\_\allowbreak{}BOUNDED\_\allowbreak{}DIAGNOSTIC & M8 and M16 are two frozen representative mixed-field tangents to which the metric framework can be applied. & Use illustrates or motivates. & Mesh-independent law, UEL revalidation, or engineering qualification. \\
C20 & T3\_\allowbreak{}BOUNDED\_\allowbreak{}DIAGNOSTIC & WP26B reproduced exact capture reconstruction but its selected-iterate subtraction identity failed the registered gate despite machine-level solve residuals. & Use motivates uncertainty-aware observability analysis. & WP26B proved catastrophic cancellation or a guard defect. \\
C21 & T4\_\allowbreak{}LIMITATION & WP27 observability labels and class counts remain diagnostic-only and are excluded from confirmatory Paper 2 conclusions. & Use as an explicit limitation and provenance note. & Use WP27 labels as validated truth or claim counts as classification accuracy. \\
C22 & T4\_\allowbreak{}LIMITATION & Local linear response metrics do not by themselves determine Newton acceptance, iteration count, basin, branch sequence, or final root. & Use does not determine by itself. & Predict nonlinear convergence or basin from eta\_\allowbreak{}z, rho\_\allowbreak{}g, or kappa. \\
C23 & T4\_\allowbreak{}LIMITATION & A tangent-only matrix perturbation is not automatically residual, mixed-FE, physical, conservation, or bound-preserving stabilization. & Use numerical guard or matrix perturbation. & Call KSTAB inf-sup stabilization or conservation correction. \\
C24 & T4\_\allowbreak{}LIMITATION & P4-r1 coverage is confined to registered dimension-2/4/8/16 SPD, nonsymmetric, and saddle-point-like decimal families and their disjoint holdout. & Use within the registered families. & Universal uncertainty model or mesh-independent coverage law. \\
C25 & T1\_\allowbreak{}THEOREM\_\allowbreak{}PLUS\_\allowbreak{}T2\_\allowbreak{}VALIDATION & A defensible assessment requires declared coordinates, separate matrix/operator/directional metrics, and uncertainty-qualified response evidence. & Use as the bounded methodological contribution. & Universal guard decision rule, optimal guard, nonlinear prediction, or engineering qualification. \\
C26 & T1\_\allowbreak{}THEOREM & A nonnegative computed response is uncertainty-qualified distinguishable from zero when d\_\allowbreak{}obs-U\textgreater{}0 under a declared truth contract and route-specific uncertainty model. & Use distinguishable from zero only with the truth contract, route, norm, and uncertainty model declared. & Observable, accurate, important, safe, harmful, or nonlinear consequence without separate evidence. \\
C27 & T2\_\allowbreak{}VALIDATED\_\allowbreak{}REGISTERED & For the frozen M8 binary-operand system, the direct-difference response is distinguishable from zero under its registered direct-route uncertainty envelope. & Use distinguishable from zero and conservative relative-error upper bound for this registered route. & Direct subtraction is generally unreliable, guard safety or harm, physical importance, nonlinear cause, M16 transfer, or mesh law. \\
C28 & T2\_\allowbreak{}VALIDATED\_\allowbreak{}REGISTERED & For the frozen M8 binary-operand system, the response-equation result is distinguishable from zero under its registered response-route uncertainty envelope. & Use distinguishable from zero under this registered response-route envelope. & Universal superiority, guard safety or harm, physical importance, nonlinear cause, M16 transfer, or mesh law. \\
C29 & T2\_\allowbreak{}VALIDATED\_\allowbreak{}REGISTERED & For this frozen M8 system, the response-equation route produced a substantially tighter registered uncertainty envelope than the direct-difference route. & Use tighter for this M8 case only. & Response equation is universally superior, optimal, cheaper, or more reliable. \\
C30 & T4\_\allowbreak{}LIMITATION & The M8 result alone does not transfer to M16; the M16 decision is supported only by its separate registered qualification, and neither case transfers to nonlinear trajectories, mesh behavior, finite-element stability, guard safety or harm, or engineering qualification. & Use independently qualified, repeated on two frozen operators, and case-specific. & Mesh law, dimension trend, population generality, nonlinear classification, safety or harm, or engineering qualification. \\
C\_\allowbreak{}M16\_\allowbreak{}01 & T2\_\allowbreak{}VALIDATED\_\allowbreak{}REGISTERED & The route-specific assessment procedure was completed on a second independently frozen coupled finite-element tangent. & Use completed on a second independently frozen tangent. & Generalized to finite-element systems, mesh trend, dimension law, or engineering qualification. \\
C\_\allowbreak{}M16\_\allowbreak{}02 & T2\_\allowbreak{}VALIDATED\_\allowbreak{}REGISTERED & The M16 direct-difference response was distinguishable from zero under its registered direct-route uncertainty envelope. & Use distinguishable from zero under this registered route envelope. & Direct-route reliability in general, safety or harm, nonlinear consequence, or mesh law. \\
C\_\allowbreak{}M16\_\allowbreak{}03 & T2\_\allowbreak{}VALIDATED\_\allowbreak{}REGISTERED & The M16 response-equation response was distinguishable from zero under its registered response-route uncertainty envelope. & Use distinguishable from zero under this registered route envelope. & Universal route superiority, safety or harm, nonlinear consequence, or mesh law. \\
C\_\allowbreak{}TWO\_\allowbreak{}OPER\allowbreak{}ATOR\_\allowbreak{}01 & T2\_\allowbreak{}VALIDATED\_\allowbreak{}REGISTERED & The same declared assessment architecture was completed for two independently frozen mixed-field tangents. & Maximum wording: repeated on two frozen operators. & Generalization, mesh convergence, dimension dependence, route ranking, safety or harm, nonlinear consequence, or engineering qualification. \\
\bottomrule
\end{xltabular}

\normalsize

\section*{Data and code availability}
The non-confidential reproducibility artifact is archived in Zenodo at \href{https://doi.org/10.5281/zenodo.21833851}{https://doi.org/10.5281/zenodo.21833851}. The public source repository is \href{https://github.com/yqma7980/scale-explicit-matrix-perturbation-observability}{https://github.com/yqma7980/scale-explicit-matrix-perturbation-observability}. The release excludes Abaqus executables, user-element source, solver state histories, and engineering project data.

\end{document}
